% =====================================================================
%  2.1 COLOR IMAGES AND LINEAR FILTERS  -  Người viết: TV2
%  CHỈ SỬA FILE NÀY. Viết tiếng Anh, tự viết, có ví dụ tính tay.
% =====================================================================
\subsection{Color images as arrays and linear color filters}\label{sec:theory-filters}
\writer{TV2}
% Gợi ý nội dung:
%  - ảnh = mảng m x n x 3; kiểu uint8, vì sao cần đổi sang số thực (Q1)
%  - pixel p = [R;G;B]; bộ lọc tuyến tính p' = M p; vì sao (M p^T)^T = p M^T
%    cho phép tính cả ảnh bằng một phép nhân ma trận
%  - ý nghĩa từng hàng / cột của M; ma trận đường chéo (Saturate), ma trận
%    hoán vị (Permute), phép chiếu (Red, DeleteGreen, Gray)
%  - phép biến đổi KHÔNG tuyến tính: đảo màu (affine), gamma (Task 9, 14)
%  - làm tròn và cắt về [0, 255] sau khi lọc
\todo{viết mục 2.1}
